%======================================================================%
\section{Yukawa from Overlap Integrals}
\label{sec:yukawa}
%======================================================================%

This section develops the programme step \textbf{D1} of
Section~\ref{sec:predictions:program}: evaluate the cubic overlap
\eqref{eq:breaking-yukawa} on $\hat Q$ with Higgs data from $\Hess V$.
What is \textbf{proved} below is the reduction of the overlap to cellular
pairing on $\mathcal{T}_7$ and the leading-order capacity-weight structure of
$y_{ij}$. What remains \textbf{conditional} is identification of
$\Phi_{\mathrm{Higgs}}$ with a specific unstable Hessian mode and the full
GUT breaking chain (Conjecture~\ref{conj:breaking-derived}).

%----------------------------------------------------------------------%
\subsection{Overlap form on the mirror quotient}
\label{sec:yukawa:overlap}
%----------------------------------------------------------------------%

Fix the stabilized Stone lift $Q$ and mirror quotient $\hat Q=Q/\sigma$
(Definition~\ref{def:Qhat}). Let $\omega_\iota\in
H^1(\hat Q,V_{16})^{\sigma=-1,\,\iota}$ denote the observer $1$-class of
Definition~\ref{def:observer}, represented by a $\sigma$-anti-invariant
Dolbeault--Serre $1$-form on $(Q,V_{16})$.

\begin{definition}[Observer overlap form]\label{def:overlap-form}
The \emph{observer overlap form} on $\hat Q$ is the antisymmetric pairing
\begin{equation}\label{eq:overlap-form}
  \langle\alpha,\beta\rangle_{\mathrm{ov}}
  \;:=\;
  \int_{\hat Q}\alpha\wedge\beta\wedge\omega_{\mathrm{vol}},
  \qquad
  \alpha,\beta\in H^1(\hat Q,V_{16})^{\sigma=-1},
\end{equation}
where $\omega_{\mathrm{vol}}$ is the $\sigma$-invariant K\"ahler volume form
induced by the $H$-invariant coset metric $K$ on the $\mathbf{16}$-bundle.
For a Higgs fluctuation $\Phi_{\mathrm{Higgs}}\in\Omega^0(\hat Q,V_{16})$
drawn from the unstable $\Hess V$ spectrum at $\orderparam_\star$
(Lemma~\ref{lem:breaking-hess}), the \emph{cubic Yukawa overlap} is
\begin{equation}\label{eq:yukawa-overlap}
  \mathcal{Y}_{ij}
  \;:=\;
  \int_{\hat Q}\omega_i\wedge\omega_j\wedge\Phi_{\mathrm{Higgs}},
  \qquad i,j\in\{1,2,3\},
\end{equation}
with $\omega_i$ the representative of $|\omega_i\rangle$. This is the
precise form of \eqref{eq:breaking-yukawa}.
\end{definition}

\begin{remark}[Epistemic split at the definition]
The overlap form \eqref{eq:overlap-form} is a standard cup-product pairing
on $H^\bullet(\hat Q,V_{16})$ restricted to the anti-invariant sector; it is
\textbf{defined}, not conjectured. The identification
$\mathcal{Y}_{ij}=y_{ij}$ in the low-energy Lagrangian is
\textbf{conditional} on Conjecture~\ref{conj:breaking-derived} and on
choosing $\Phi_{\mathrm{Higgs}}$ in the $\mathbf{10}_H$ (or $\mathbf{126}_H$)
channel of the breaking chain.
\end{remark}

%----------------------------------------------------------------------%
\subsection{Cellular reduction on $\mathcal{T}_7$}
\label{sec:yukawa:cellular}
%----------------------------------------------------------------------%

On the stabilized Albert tower $\mathcal{T}_7$ at $\Mirror^7(\Void)$, write
$C_1,C_2,C_3$ for the three $\mathbf{16}$-cells attached to diagonal
directions $e_1,e_2,e_3$ of $\mathfrak{d}$ (Lemma~\ref{lem:cellular}).
Let $\mathcal{F}_{16}$ be the cellular cosheaf of Appendix
\ref{sec:appendix:quotient}.

\begin{lemma}[Cubic overlap as cellular pairing]\label{lem:yukawa-cellular}
Let $\omega_i$ be the cellular cocycle representative of
$|\omega_i\rangle\in H^1(\mathcal{T}_7,\mathcal{F}_{16})^{\sigma=-1,\,i}$.
Then the cubic overlap \eqref{eq:yukawa-overlap} reduces to a finite sum
over $\mathcal{T}_7$:
\begin{equation}\label{eq:yukawa-cellular}
  \mathcal{Y}_{ij}
  \;=\;
  \sum_{k=1}^3
  \mathrm{pair}\!\bigl(C_i,C_j,C_k;\Phi_{\mathrm{Higgs}}\bigr),
\end{equation}
where $\mathrm{pair}(C_i,C_j,C_k;\Phi)$ is the oriented triple intersection
number of the three $1$-cells at the $\mathbf{10}$-hub, weighted by the
$\Phi$-value on the shared $0$-simplex. Under the Stone lift
(Proposition~\ref{prop:stone-lift}),
$H^1(\mathcal{T}_7,\mathcal{F}_{16})^{\sigma=-1}\cong
H^1(\hat Q,V_{16})^{\sigma=-1}$, and \eqref{eq:yukawa-cellular} agrees with
the continuum integral \eqref{eq:yukawa-overlap}.
\end{lemma}

\begin{proof}
On $\mathcal{T}_7$, $1$-cochains are triples $(\phi_1,\phi_2,\phi_3)$ with
cocycle constraint $\phi_1+\phi_2+\phi_3=0$ at the hub $H_{10}$
(Lemma~\ref{lem:cellular}, \eqref{eq:cell-cocycle}). The quotient
\eqref{eq:cell-quotient} identifies
$H^1(\mathcal{T}_7,\mathcal{F}_{16})^{\sigma=-1}$ with three
$e_i$-weight classes; Proposition~\ref{prop:triality} fixes their orientation.
The wedge product of two observer $1$-forms and a $0$-form localizes to the hub
$0$-simplex where all three cells meet: only terms with distinct cell indices
contribute, and $\sigma$-anti-invariance forces
$\mathrm{pair}(C_i,C_j,C_k)=0$ unless $\{i,j,k\}=\{1,2,3\}$ with orientation
$\epsilon_{ijk}$. Proposition~\ref{prop:stone-lift} identifies cellular
cohomology with $H^1(\hat Q,V_{16})^{\sigma=-1}$, transporting the sum to the
integral \eqref{eq:yukawa-overlap}.
\end{proof}

\begin{corollary}[Diagonal dominance at tree level]\label{cor:yukawa-diagonal}
At the EIII vacuum $\orderparam_\star$, tree-level $\Hess V=0$ on
$\mathfrak{m}$ (Lemma~\ref{lem:breaking-hess}). If
$\Phi_{\mathrm{Higgs}}$ is the leading unstable mode localized to cell
$\iota$ after Coleman--Weinberg lifting, then
$\mathcal{Y}_{ij}\neq0$ only when $i=j=\iota$ at tree level; off-diagonal
entries arise at higher order in the radiative expansion or from sequential
$SU(3)_F$ breaking (Conjecture~\ref{conj:flavon}).
\end{corollary}

%----------------------------------------------------------------------%
\subsection{Capacity-weight structure}
\label{sec:yukawa:weights}
%----------------------------------------------------------------------%

Descend the sector masses $(1/27,10/27,16/27)$ of Theorem~\ref{thm:partition}
through the above--below correspondence (Proposition~\ref{prop:above-below})
to observer-slot weights
\begin{equation}\label{eq:yukawa-weights}
  w_1:w_2:w_3 = 1:10:16,
  \qquad
  w_\iota := \frac{m_\iota}{1/27},
\end{equation}
with $(m_1,m_{10},m_{16})=(1/27,10/27,16/27)$ the capacity targets of
Theorem~\ref{thm:partition} on $\mathbf{27}=\mathbf{1}\oplus\mathbf{10}\oplus\mathbf{16}$.

\begin{proposition}[Leading Yukawa factorization]\label{prop:yukawa-factor}
Assume Lemma~\ref{lem:yukawa-cellular} and that $\Phi_{\mathrm{Higgs}}$ is
$H$-invariant with VEV $\langle\Phi_{\mathrm{Higgs}}\rangle$ at the
$\iota$-restricted breaking minimum. Then the cubic overlap factorizes at
leading order as
\begin{equation}\label{eq:yukawa-factor}
  \mathcal{Y}_{ij}
  \;=\;
  y_0\,\sqrt{w_i\,w_j}\,\mathcal{P}_{ij}
  \;+\;\mathcal{O}(\text{radiative}),
\end{equation}
where $y_0$ is a dimensionless coupling set by the unstable-mode residue of
$\Hess V$ and $\mathcal{P}_{ij}\in\{0,1\}$ encodes the oriented hub pairing
of Lemma~\ref{lem:yukawa-cellular}. In the diagonal-dominant regime of
Corollary~\ref{cor:yukawa-diagonal}, $\mathcal{P}_{ij}=\delta_{ij}$.
\end{proposition}

\begin{proof}[Proof sketch]
Each observer class $|\omega_\iota\rangle$ is normalized on cell $C_\iota$ with
weight $\|{\omega_\iota}\|^2\propto w_\iota$ by capacity balance on
$\mathcal{G}_{\mathrm{Albert}}$ (Definition~\ref{def:lifetime-index},
Proposition~\ref{prop:capacity-lifetime-dual}). The hub pairing is bilinear
in the two $1$-forms and linear in $\Phi_{\mathrm{Higgs}}$; leading-order
multiplicativity of capacity weights under $\mathrm{Bal}_{\mathcal{C}}$ on
the $1|10|16$ partition yields the $\sqrt{w_i w_j}$ structure. The square
root arises because $1$-forms carry capacity weight $\sqrt{w_\iota}$ while
masses scale as $w_\iota$ at fixed Higgs VEV
(Proposition~\ref{prop:gen-ordering}).
\end{proof}

\begin{remark}[Theorem vs.\ conjecture]
Lemma~\ref{lem:yukawa-cellular} is a \textbf{theorem} (cellular cohomology
plus Stone lift). Proposition~\ref{prop:yukawa-factor} is a
\textbf{proposition} contingent on the Higgs-mode identification of
Conjecture~\ref{conj:breaking-derived}; the $\sqrt{w_i w_j}$ scaling itself is
\textbf{leading-order} capacity descent, not yet a sub-percent fit.
\end{remark}

%----------------------------------------------------------------------%
\subsection{Mass matrix at $\orderparam_\star$}
\label{sec:yukawa:mass-matrix}
%----------------------------------------------------------------------%

\begin{theorem}[Leading mass matrix, conditional]\label{thm:yukawa-mass}
Assume Conjecture~\ref{conj:breaking-derived}, Proposition~\ref{prop:yukawa-factor},
and a common electroweak Higgs VEV $v$. Define the fermion mass matrix by
\begin{equation}\label{eq:mass-matrix-def}
  M_{ij}
  \;:=\;
  v\,\mathcal{Y}_{ij}
  \;=\;
  v\,y_0\,\sqrt{w_i\,w_j}\,\mathcal{P}_{ij}
  \;+\;\mathcal{O}(\text{radiative}).
\end{equation}
At the EIII reference $\orderparam_\star$ with diagonal-dominant hub pairing
($\mathcal{P}_{ij}=\delta_{ij}$), $M$ is positive semidefinite and
eigenvalues satisfy
\begin{equation}\label{eq:mass-eigenvalues}
  m_1:m_2:m_3 \;=\; w_1:w_2:w_3 \;=\; 1:10:16
\end{equation}
at leading order. The theorem is \textbf{conditional}: it does not prove the
breaking chain or the Higgs quantum numbers; given those inputs, the mass
hierarchy follows from capacity weights already fixed by
Theorem~\ref{thm:partition}.
\end{theorem}

\begin{proof}[Proof sketch]
With $\mathcal{P}_{ij}=\delta_{ij}$, $M=\mathrm{diag}(v y_0\sqrt{w_1},
v y_0\sqrt{w_2}, v y_0\sqrt{w_3})$. Squaring,
$M_{ii}^2\propto w_i$, hence eigenmass ordering $m_3>m_2>m_1$. The ratio
$m_3:m_2:m_1=w_3:w_2:w_1=16:10:1$ is exact at this order.
\end{proof}

\begin{equation}\label{eq:mass-matrix-explicit}
  M
  \;=\;
  v\,y_0
  \begin{pmatrix}
    \sqrt{w_1} & 0 & 0 \\
    0 & \sqrt{w_2} & 0 \\
    0 & 0 & \sqrt{w_3}
  \end{pmatrix}
  \;=\;
  v\,y_0
  \begin{pmatrix}
    1 & 0 & 0 \\
    0 & \sqrt{10} & 0 \\
    0 & 0 & 4
  \end{pmatrix},
\end{equation}
using $w_1:w_2:w_3=1:10:16$ so $\sqrt{w_1}:\sqrt{w_2}:\sqrt{w_3}=1:\sqrt{10}:4$.
Off-diagonal entries are $\mathcal{O}(\text{flavon})$ once sequential
$SU(3)_F$ breaking activates (Conjecture~\ref{conj:flavon}); the full
symmetric leading ansatz is
\begin{equation}\label{eq:mass-matrix-full}
  M_{ij} \;=\; v\,y_0\,\sqrt{w_i\,w_j}\,\mathcal{P}_{ij},
  \qquad
  \mathcal{P}_{ij}\in\{0,1\}.
\end{equation}

%----------------------------------------------------------------------%
\subsection{Numerical leading spectrum}
\label{sec:yukawa:numerical}
%----------------------------------------------------------------------%

Table~\ref{tab:yukawa-leading} records the leading-order spectrum implied by
Theorem~\ref{thm:yukawa-mass}. Slot $\iota=2$ carries capacity weight
$w_2=10/27$; under the sector map of Prediction~\ref{pred:sector-ratios},
this is the \textbf{muon} slot: $m_\mu/m_e\sim\sqrt{w_2/w_1}=\sqrt{10}\approx
3.2$ at the Yukawa level and $m_\mu/m_e\sim w_2/w_1=10$ at the mass level
(Proposition~\ref{prop:gen-ordering}).

\begin{table}[ht]
\centering
\small
\renewcommand{\arraystretch}{1.12}
\begin{tabular}{@{}cccccc@{}}
\toprule
\textbf{Slot $\iota$} & $w_\iota$ & $\sqrt{w_\iota}$ (rel.) &
$\mathcal{Y}_{\iota\iota}$ (rel.) & $m_\iota$ (rel.) &
\textbf{Leading identification} \\
\midrule
$1$ & $1/27$  & $1$        & $1$        & $1$  & lightest (e.g.\ $e$) \\
$2$ & $10/27$ & $\sqrt{10}$ & $\sqrt{10}$ & $10$ & muon slot \\
$3$ & $16/27$ & $4$        & $4$        & $16$ & heaviest (e.g.\ $\tau$) \\
\bottomrule
\end{tabular}
\caption{Leading-order Yukawa and mass hierarchy from overlap factorization
\eqref{eq:yukawa-factor} at $\orderparam_\star$. Absolute scales require
step~\textbf{D3} (scale identification). Ratios are exact at this order;
RG and flavon corrections are programme step~\textbf{D2}.}
\label{tab:yukawa-leading}
\end{table}

\begin{corollary}[D1 contact with predictions]\label{cor:yukawa-d1}
Section~\ref{sec:yukawa} \textbf{partially closes} development step~\textbf{D1}
of Section~\ref{sec:predictions:program}: the cubic overlap
\eqref{eq:breaking-yukawa} is defined on $\hat Q$
(Definition~\ref{def:overlap-form}), reduced to $\mathcal{T}_7$
(Lemma~\ref{lem:yukawa-cellular}), and evaluated at leading order
(Proposition~\ref{prop:yukawa-factor}, Theorem~\ref{thm:yukawa-mass}).
Precision match to PDG masses and to JUNO $\Delta m^2$ ratios remains
\textbf{conditional} on steps~\textbf{D2}--\textbf{D3}.
\end{corollary}

\begin{conjecture}[Full Yukawa sector]\label{conj:yukawa-full}
Radiatively lifted $\Hess V$ on the $\mathbf{27}$-sector, filtered through
$\Pi_{\sigma=-1}$, selects $\Phi_{\mathrm{Higgs}}\in\mathbf{10}_H$ with
VEV $v$, and generates the complete $3\times3$ Yukawa matrices for quarks
and leptons (including CKM/PMNS from off-diagonal $\mathcal{P}_{ij}$) with
no additional Yukawa postulate. This upgrades Conjecture~\ref{conj:breaking-derived}
from a schematic overlap to a calculable mass spectrum.
\end{conjecture}